\section*{Introduction}
  Ce rapport présente les découvertes effectuées lors de mon stage de première année à l'\enseirb dans le cadre de l'animation d'un atelier de découverte destiné à des élèves de seconde. Très intéressé par l'enseignement et le partage de connaissances, j'ai choisi de postuler à l'offre de stage proposée par l'\enseirb en mars 2026, dont l'objectif était le pilotage de plusieurs ateliers de vulgarisation à destination des trente-six élèves de classe de seconde qui viennent visiter l'école dans le cadre de leur stage. Ce pilotage s'est fait en binôme avec un.e autre étudiant.e de première année de l'école et a concerné la thématique "\textit{Apprendre à discuter avec l'intelligence artificielle}". La présentation de l'atelier se fait à destination de binômes d'élèves de seconde, sur une durée de trois heures.\par 
  L'\enseirb est une école d'ingénieurs située à Talence, en France, qui propose des formations dans les domaines de l'informatique, de l'électronique et des télécommunications. L'école est reconnue pour son excellence académique et son engagement envers l'innovation technologique notamment du point de vue associatif. \par
J'ai personnellement été intéressé par l'opportunité de participer à cet atelier, très enthousiaste à l’idée de partager mes connaissances et de contribuer à démystifier les idées reçues, les doutes et les fantasmes liés à l’intelligence artificielle et contribuer à un usage plus sain de tels outils, notamment chez les plus jeunes, dans un cadre scolaire. \par
  Ce stage s’est articulé autour de deux phases distinctes : la conception de l’atelier, puis son animation. à trente-six élèves de seconde. \\\par
  Dans ce rapport nous aborderons d'un aspect chronologique les étapes de conceptualisation et de présentation de l'atelier, ainsi que mes retours d'expérience sur l'animation de celui-ci. 
